% =====================================================================
%   §6. Счётчики с короткой памятью HyperLogLog
% =====================================================================
\section{Счётчики с короткой памятью HyperLogLog и эксперимент Стенли Мильграма}
\label{sec:hll}

В этом, заключительном параграфе главы мы рассмотрим ещё один сюжет на стыке теории чисел, теории вероятностей и информатики --- задачу, в которой важно \emph{подсчитывать число различных объектов} в огромном потоке данных, не имея возможности хранить их все. Это пригодится для задачи, на первый взгляд несвязанной с криптографией: измерения «социальных расстояний» в больших сетях, и проверки знаменитой гипотезы Стенли Мильграма о шести рукопожатиях.

Изложение в этом параграфе опирается на главу 7 книги Литвак и Райгородского «Кому нужна математика?»\footnote{Литвак~Н.\,В., Райгородский~А.\,М. \emph{Кому нужна математика? Содержательное введение в математику и computer science}. М.: Манн, Иванов и Фербер, 2017. Гл.~6, 7 и приложения.}, которую мы рекомендуем для самостоятельного изучения.

% ---------------------------------------------------------------------
\subsection{Задача подсчёта различных элементов}

Представьте, что вы работаете в большой интернет-компании и хотите узнать: сколько \emph{уникальных} пользователей зашло сегодня на сайт? Не общее число обращений (его легко посчитать инкрементируя счётчик при каждом запросе), а именно \emph{различных} людей.

\paragraph*{Точное решение.} Простое и очевидное: завести множество (например, хеш-таблицу из §\ref{sec:hashing}), и при каждом приходящем идентификаторе пользователя проверять, есть ли он там; если нет --- добавить и увеличить счётчик. По окончании дня размер множества --- ответ.

\paragraph*{Проблема.} Если число уникальных пользователей --- сотни миллионов или миллиарды, такая таблица занимает гигабайты памяти. Для одной задачи на одном сервере --- допустимо. А если таких задач тысячи, и каждая на каждом из тысяч серверов? Память становится узким местом.

Возникает естественный вопрос: можно ли подсчитать число различных элементов в потоке \emph{приближённо}, но используя \emph{гораздо меньше памяти}? Например, дать ответ с относительной погрешностью $\pm 2\%$, заняв всего несколько килобайт?

Удивительно, но ответ положительный. И построение такого алгоритма опирается на красивые идеи теории вероятностей и --- что нам особенно интересно --- использует хеш-функции, родственные тем, что мы изучали в §\ref{sec:hashing}.

% ---------------------------------------------------------------------
\subsection{Идея: «самый длинный хвост нулей»}

Алгоритм, который мы сейчас опишем (он называется \term{HyperLogLog}, предложен Филиппом Флажоле и его коллегами в 2007 году), основан на одной красивой вероятностной идее.

Пусть у нас есть «хорошая» хеш-функция $h$, которая для каждого входного объекта (например, идентификатора пользователя) выдаёт случайную на вид последовательность из $L$ битов. Слово «случайную» означает: для разных входов биты выглядят как независимые подбрасывания монеты\footnote{Конечно, $h$ детерминированна, но её устройство таково, что заранее, не считая, угадать биты невозможно. На практике используются криптографические или близкие к ним хеш-функции, обладающие именно таким свойством «равномерности».}.

\paragraph*{Наблюдение.} Если число $h(x)$ действительно ведёт себя случайно, то рассмотрим, сколько \emph{нулей подряд} стоит в начале двоичной записи $h(x)$.
\begin{itemize}
  \item Вероятность, что первый бит --- ноль: $1/2$.
  \item Вероятность, что первые два бита --- нули: $1/4$.
  \item Вероятность, что первые $k$ битов --- нули: $1/2^k$.
\end{itemize}

Наглядно эта закономерность показана на рис.~\ref{fig:p6-1}.

\begin{figure}[H]
\centering
\begin{tikzpicture}[every node/.style={font=\small\ttfamily}]
  \node at (0,0)   {010101...};
  \node at (0,-0.5){0000111...};
  \node at (0,-1)  {001011...};
  \node at (0,-1.5){\textcolor{prosvRed}{00000}10...};
  \node at (0,-2)  {1110...};
  \node[font=\small\rmfamily, prosvBlue] at (4, -1.5) {длинный «хвост нулей» --- редкое событие};
  \draw[->, prosvRed] (2.5, -1.5) -- (0.8, -1.5);
\end{tikzpicture}
\caption{Иллюстрация двоичных хешей. Чем длиннее «хвост нулей» в начале, тем реже такое событие.}
\label{fig:p6-1}
\end{figure}

Иначе говоря, если в потоке встречается элемент $x$, такой что $h(x)$ начинается, скажем, с $20$ нулей подряд, это значит, что в потоке \emph{скорее всего побывало порядка $2^{20} \approx 10^6$ различных элементов} --- иначе крайне маловероятно увидеть такое редкое событие.

\paragraph*{Базовый алгоритм («одиночный счётчик»).}

\begin{algorithmbox}[title={Алгоритм: наивный счётчик различных элементов}]
\begin{enumerate}
  \item Заведём одну переменную $M = 0$.
  \item Для каждого приходящего элемента $x$ потока:
  \begin{enumerate}
    \item вычислим $h(x)$;
    \item посчитаем число $\rho(h(x))$ --- длину начального блока нулей в $h(x)$, плюс единица;
    \item обновим $M \leftarrow \max(M, \rho(h(x)))$.
  \end{enumerate}
  \item На выходе оценим число различных элементов как $\widehat{n} = 2^M$.
\end{enumerate}
\end{algorithmbox}

\paragraph*{Память.} Чтобы хранить $M$ от $0$ до $L$ (где $L$ --- длина хеша, скажем, 32 бита), достаточно \emph{шести битов памяти}. Шесть битов --- чтобы оценить число различных элементов в потоке любой длины!

\paragraph*{Проблема.} Эта оценка очень шумная: одно «случайно длинное» обнуление сбивает её в большую сторону, а его отсутствие --- в меньшую. Дисперсия чудовищная.

% ---------------------------------------------------------------------
\subsection{Идея «несколько счётчиков»}

Чтобы сгладить случайные колебания, разделим поток на много подпотоков и усредним результаты. Это и есть алгоритм HyperLogLog.

\begin{algorithmbox}[title={Алгоритм: HyperLogLog (упрощённая схема)}]
\begin{enumerate}
  \item Выберем число $b$ --- скажем, $b = 10$, --- и заведём $2^b = 1024$ независимых счётчика $M_0, M_1, \dots, M_{2^b - 1}$. Каждый счётчик --- 6 битов.
  \item Для каждого приходящего элемента $x$:
  \begin{enumerate}
    \item вычислим $h(x)$ длины $L$;
    \item первые $b$ битов $h(x)$ интерпретируем как номер счётчика $j$, $0 \le j < 2^b$;
    \item к остальным $L - b$ битам применяем правило: $\rho = $ длина начального блока нулей плюс 1;
    \item обновляем: $M_j \leftarrow \max(M_j, \rho)$.
  \end{enumerate}
  \item На выходе оцениваем
  \[
    \widehat{n} = \alpha_b \cdot 2^b \cdot \left( \frac{1}{2^b} \sum_{j=0}^{2^b-1} 2^{-M_j} \right)^{-1},
  \]
  где $\alpha_b$ --- известная константа корректировки (для $b = 10$ примерно $0.7213$).
\end{enumerate}
\end{algorithmbox}

\paragraph*{Почему так?} Формула в третьем шаге --- это \emph{гармоническое среднее} величин $2^{M_j}$, а не простое арифметическое. Гармоническое среднее устойчиво к выбросам (одно случайно большое $M_j$ его существенно не искажает) --- именно поэтому HyperLogLog так точен.

\begin{theorem}{точность HyperLogLog}{thm:p6-1}
Относительная погрешность оценки $\widehat{n}$ по сравнению с истинным значением $n$ имеет порядок $1.04/\sqrt{2^b}$. При $b = 10$ это примерно $3.25\%$; при $b = 14$ --- около $0.81\%$.
\end{theorem}

\paragraph*{Память.} При $b = 14$ имеется $2^{14} = 16\,384$ счётчиков по 6 битов --- всего $\approx 12$ КБ. И этим объёмом мы оцениваем число уникальных элементов в потоке длиной до $\approx 10^{18}$ с точностью около $1\%$! (рис.~\ref{fig:p6-2})

\begin{figure}[H]
\centering
\begin{tikzpicture}[
  every node/.style={font=\footnotesize},
  cell/.style={draw=prosvBlue, fill=prosvLight, minimum width=0.7cm, minimum height=0.7cm},
  hl/.style={draw=prosvRed, fill=prosvCream, minimum width=0.7cm, minimum height=0.7cm, thick}
]
  \node[font=\small, prosvBlue] at (-2, 0.7) {Счётчики:};
  \node[cell] at (-1.3, 0) {3};
  \node[cell] at (-0.6, 0) {1};
  \node[hl]   at (0.1, 0)  {7};
  \node[cell] at (0.8, 0)  {2};
  \node[cell] at (1.5, 0)  {4};
  \node[cell] at (2.2, 0)  {0};
  \node[cell] at (2.9, 0)  {3};
  \node[cell] at (3.6, 0)  {1};
  \node[cell] at (4.3, 0)  {2};
  \node[cell] at (5.0, 0)  {5};
  \node at (5.7, 0)   {$\dots$};

  \node[align=left, font=\footnotesize, prosvRed] at (8, 0) {\textit{редкий «хвост»} ---\\ скорее всего \\ много элементов};
  \draw[->, prosvRed] (8, -0.5) -- (0.4, -0.3);

  \node[font=\footnotesize\itshape, prosvBlue, align=center] at (3, -1.2)
    {усреднение по всем счётчикам с помощью гармонического среднего};
\end{tikzpicture}
\caption{Множество счётчиков HyperLogLog. Каждый запоминает максимальную длину «хвоста нулей» для своего подпотока.}
\label{fig:p6-2}
\end{figure}

\begin{interestbox}[title={Это интересно: где работает HyperLogLog}]
HyperLogLog встроен в Redis (популярную систему кэширования), в Google BigQuery (язык запросов к огромным базам данных), в систему аналитики ВКонтакте, в инфраструктуру PostgreSQL и многих других сервисов. Когда вы видите в личном кабинете рекламной системы оценку «охват: 3{,}5\,млн уникальных пользователей», очень вероятно, что под капотом работает именно этот алгоритм.
\end{interestbox}

% ---------------------------------------------------------------------
\subsection{Эксперимент Стенли Мильграма}
\label{subsec:milgram}

Теперь, когда у нас в руках мощный инструмент для подсчёта уникальных объектов, обратимся к классической задаче из социологии, которую Литвак и Райгородский подробно разбирают в главе 7 своей книги.

\paragraph*{Историческое введение.} В 1960-х годах американский социальный психолог \term{Стенли Мильграм} (1933--1984) задался необычным вопросом: \emph{насколько на самом деле «велик» мир}? Точнее: если выбрать наугад двух человек на разных концах страны, через сколько посредников можно их связать?

В 1967 году Мильграм провёл свой знаменитый эксперимент. Участникам в Канзасе и Небраске были даны письма с просьбой передать их некоему адресату в Бостоне, причём \emph{нельзя} было отправлять напрямую --- только через личного знакомого, который, как кажется участнику, может быть ближе к получателю; тот, в свою очередь, выбирал следующего знакомого, и так далее.

\paragraph*{Удивительный результат.} Большинство писем, дошедших до адресата, преодолевали путь \emph{в среднем за 5--6 пересылок}. Так родилась знаменитая гипотеза «шести рукопожатий»: между любыми двумя людьми на Земле в среднем стоит лишь шесть посредников.

\begin{interestbox}[title={Это интересно: «шесть рукопожатий» в литературе}]
Идея «малого мира» проникла в массовую культуру: пьеса Джона Гуэра «Six Degrees of Separation» (1990), фильм с тем же названием, многочисленные «эстафеты знакомств». В науке термин \term{small-world network} стал важным понятием в теории графов и сетевой социологии.
\end{interestbox}

\paragraph*{Современная проверка.} Эксперимент Мильграма страдал серьёзными методологическими дефектами: 1/\,подавляющее большинство писем (около 80\%) до адресата так и не дошли; 2/\,выборка была очень ограниченной (всего около 300 писем); 3/\,«знакомые» отбирались не случайно, а с предположением о близости к получателю.

В современную эпоху, когда подавляющая часть знакомств и контактов перешла в цифровой формат, задачу можно проверить иначе: \emph{построить граф знакомств целиком} и измерить кратчайшие расстояния. Именно это в 2011 году проделали исследователи в рамках одной известной социальной сети --- они проанализировали данные \emph{всех} её пользователей (на тот момент --- около 720 миллионов человек) и более 69 миллиардов «дружеских связей».

\paragraph*{Что они обнаружили.} Средняя длина кратчайшего пути между двумя случайно выбранными пользователями оказалась равной \emph{4{,}74 рукопожатия}. Иначе говоря, в современных социальных сетях гипотеза Мильграма не просто подтверждается, а даже \emph{улучшается}: между двумя пользователями стоит в среднем менее пяти промежуточных знакомых, а не шесть. Принято округлять это число до 4 --- отсюда выражение «четыре рукопожатия» (рис.~\ref{fig:p6-3}).

\begin{figure}[H]
\centering
\begin{tikzpicture}[
  >=Stealth,
  person/.style={circle, draw=prosvBlue, fill=prosvLight, minimum size=0.7cm, font=\footnotesize, inner sep=0},
]
  \node[person] (a) at (0,0)   {А};
  \node[person] (b) at (1.8,0.8) {?};
  \node[person] (c) at (3.6,0)   {?};
  \node[person] (d) at (5.4,0.8) {?};
  \node[person] (e) at (7.2, 0)  {Я};

  \draw[thick, prosvRed] (a) -- (b) -- (c) -- (d) -- (e);

  \node[font=\footnotesize, prosvBlue, align=center] at (3.6, -1.2)
    {4 рукопожатия = 4 ребра в графе знакомств};
\end{tikzpicture}
\caption{«Четыре рукопожатия»: между Алисой и Яковом стоят всего трое промежуточных знакомых.}
\label{fig:p6-3}
\end{figure}

% ---------------------------------------------------------------------
\subsection{При чём здесь HyperLogLog?}

Возникает вопрос: \emph{как именно} вычислить среднюю длину кратчайшего пути в графе из 720 миллионов вершин и 69 миллиардов рёбер? Прямой подсчёт всех пар (а их $\binom{720{\cdot}10^6}{2} \approx 2{,}6 \cdot 10^{17}$) непосильно даже для самых мощных кластеров.

Здесь и приходит на помощь HyperLogLog (а также его «родственники» --- алгоритмы вероятностного подсчёта). Идея, идущая в реальном анализе, такова:

\begin{algorithmbox}[title={Алгоритм: оценка расстояний в графе через подсчёт «соседей»}]
\begin{enumerate}
  \item Для каждой вершины $v$ заведём HyperLogLog-счётчик $C_v$.
  \item Вначале добавим в $C_v$ только саму вершину $v$. То есть $|C_v| = 1$.
  \item На шаге $t = 1, 2, \dots, T$:
  \begin{itemize}
    \item для каждой вершины $v$ обновим $C_v \leftarrow C_v \cup \bigcup_{u: (u,v) \text{ -- ребро}} C_u^{(t-1)}$,
    \item то есть в $C_v$ добавляются все вершины, до которых дошли соседи $v$ на предыдущем шаге.
  \end{itemize}
  \item Величина $|C_v|$ после $t$ шагов --- это число вершин, до которых от $v$ существует путь длины $\le t$.
\end{enumerate}
\end{algorithmbox}

\paragraph*{Где экономия памяти.} Если бы мы хранили $C_v$ как точное множество, для одной вершины пришлось бы выделить до 720 миллионов битов памяти --- 90 мегабайт. Умножим на 720 миллионов вершин --- получим 65 петабайт. Невозможно.

А поскольку HyperLogLog-счётчик занимает 12 КБ на вершину, всё хранилище умещается в $720{\cdot}10^6 \cdot 12$ КБ $\approx 8{,}6$ ТБ --- объём, который уже умещается в один большой серверный кластер.

\paragraph*{Объединение HyperLogLog-счётчиков.} У алгоритма есть ещё одно прекрасное свойство: \term{HLL-счётчики можно объединять}. Если $C_1$ оценивает мощность множества $A$, а $C_2$ --- мощность $B$, то новый счётчик $C$, у которого $C[j] = \max(C_1[j], C_2[j])$, оценивает мощность $A \cup B$. Это позволяет распределять вычисления по тысячам машин и в конце «склеивать» результаты.

\begin{importantbox}[title={Важно: краткое резюме}]
Современное измерение «социального расстояния» в больших сетях --- это композиция нескольких идей:
\begin{enumerate}
  \item теория графов даёт постановку задачи (кратчайшие расстояния);
  \item теория чисел и хеширование (§\ref{sec:hashing}) обеспечивают «равномерное» распределение элементов;
  \item теория вероятностей и алгоритм HyperLogLog позволяют оценивать мощности множеств за десятки тысяч раз меньшую память;
  \item распределённые вычисления позволяют разнести работу по тысячам компьютеров.
\end{enumerate}
\end{importantbox}

% ---------------------------------------------------------------------
\subsection{От социологии --- к стратегическим выводам}

Эксперимент Мильграма и его цифровое продолжение --- не просто забавный факт. Из него следуют важные вещи:

\begin{itemize}
  \item \emph{Информация распространяется быстро.} Новость, попавшая в социальную сеть, при равномерной передаче «через знакомых» теоретически охватывает всю сеть за 4--5 шагов. На практике это и наблюдается с вирусными мемами.
  \item \emph{Эпидемии распространяются быстро.} В аналогичной модели заражения болезнь, передающаяся при контакте, охватывает большую долю населения за считанные «волны» контактов.
  \item \emph{Цифровая приватность хрупка.} Если ваш друг знает всё о вас, его друг --- через пять рукопожатий --- статистически связан с миллионами людей. Любая утечка приватной информации распространяется по сети быстрее, чем кажется.
  \item \emph{Поиск работы и идей.} Знаменитая работа социолога Марка Грановеттера 1973 года показала, что новые контакты, идеи, рабочие места часто приходят к нам не от близких друзей, а от «слабых связей» --- знакомых второго и третьего уровня. Малый диаметр сети объясняет, почему такая стратегия работает.
\end{itemize}

\begin{interestbox}[title={Это интересно: модель «малого мира» Уоттса--Строгаца}]
Математическая модель социальной сети, объясняющая эффект малого мира, была построена Дунканом Уоттсом и Стивеном Строгацем в 1998 году. Они показали: если у регулярной решётки добавить совсем немного случайных «дальних» рёбер, диаметр графа резко падает с линейного до логарифмического. Социальная сеть устроена именно так: большинство ваших связей --- локальные (одноклассники, коллеги), но достаточно нескольких «дальних» --- например, родственника в другом городе или знакомого по интернету, --- чтобы вся сеть оказалась «компактной».
\end{interestbox}

\begin{tasksbox}
\begin{enumerate}
  \item Объясните, почему в наивном алгоритме одиночного счётчика оценка $\widehat{n} = 2^M$ имеет такую большую дисперсию.
  \item Пусть хеш-функция выдаёт строки по 8 битов. Сколько (в среднем) различных элементов нужно подать на вход, чтобы хотя бы у одного оказалось 6 нулей в начале хеша?
  \item Объясните, почему гармоническое среднее устойчивее к выбросам, чем арифметическое. Приведите пример набора чисел, в котором эти средние существенно различаются.
  \item Оцените, сколько памяти займёт HyperLogLog при $b = 12$, и какова при этом ожидаемая относительная погрешность.
  \item* Почему в алгоритме оценки расстояний в графе нельзя просто хранить $|C_v|$ как число (это же 4 байта)? Зачем хранить целый HLL-счётчик?
  \item* Среднее число рукопожатий в графе оценено как $4{,}74$. Что в этой формулировке играет роль случайной величины, а что --- константы? Какой стандартной операции теории вероятностей соответствует «среднее число»?
\end{enumerate}
\end{tasksbox}
